% XIII LAW3M 2026 abstract
% One-page abstract prepared following the conference template.
\documentclass[12pt,a4paper]{article}
\usepackage[utf8]{inputenc}
\usepackage[T1]{fontenc}
\usepackage[english]{babel}
\usepackage{mathptmx}
\usepackage{amsmath}
\usepackage{geometry}
\usepackage{setspace}
\usepackage[normalem]{ulem}
\usepackage{microtype}
\geometry{top=2.0cm,bottom=2.0cm,left=2.2cm,right=2.2cm}
\pagestyle{empty}
\setlength{\parindent}{0pt}
\setlength{\parskip}{0pt}
\singlespacing
\sloppy
\begin{document}
\begin{center}
{\bfseries\fontsize{14}{16}\selectfont XIII LAW3M - Abstract Submission Template}\par
\vspace{6pt}
{\bfseries\fontsize{14}{16}\selectfont simulated dipolar compass-needle arrays as a model for damping-controlled magnetic avalanches}\par
\vspace{10pt}
{\bfseries\fontsize{12}{14}\selectfont \uline{Jo\~ao~P.~Sinnecker}}\par
\vspace{6pt}
{\itshape\fontsize{12}{14}\selectfont Centro Brasileiro de Pesquisas F\'isicas (CBPF), Rio de Janeiro, Brazil}\par
\end{center}
\vspace{10pt}
Artificial spin-ice arrays show how simple magnetic units coupled mainly by dipolar fields can generate collective states such as domains, hysteresis, and avalanche-like reversal~[1]. In nanolithographic systems, however, island shape, spacing, disorder, switching barriers, and damping are strongly entangled. This makes it difficult to isolate which parameter controls the observed magnetic dynamics. Here we use numerical simulations of two-dimensional arrays of compass needles as a minimal model for this problem. Each needle is represented as a magnetic dipole fixed at a lattice site and free to rotate in the plane under the total dipolar field of the other needles, an applied magnetic field, and viscous damping. The model keeps the long-range and anisotropic character of dipolar coupling, while allowing the damping to be varied independently. The central control parameter is the quality factor $Q$, which compares inertial motion with viscous damping. Low-$Q$ arrays reverse slowly and almost monotonically, whereas high-$Q$ arrays show overshoot, oscillatory relaxation, and chain-reaction switching events. We use this tunability to ask whether the dynamical regime set by damping changes the hysteresis loop and the statistics of magnetic avalanches in classical dipolar lattices. Simulations are performed for square, triangular, and honeycomb arrays, allowing the role of lattice geometry and magnetic frustration to be compared under the same physical assumptions. For each geometry and damping condition, five-ramp hysteresis cycles are calculated and followed by zero-field relaxation. The analysis extracts loop area, coercive field, remanence, avalanche-size distributions, waiting times between avalanches, and final domain statistics. Preliminary results indicate that damping is not a passive numerical parameter: it changes both the shape of the hysteresis loop and the distribution of reversal events. In particular, underdamped arrays tend to produce larger correlated switching cascades, while overdamped arrays display more gradual reversal. The triangular geometry shows additional sensitivity of the domain structure to $Q$, suggesting an interplay between frustration and inertia. The compass-needle system is therefore not proposed as a literal substitute for nanomagnetic islands, but as a transparent dipolar model in which geometry, spacing, damping, and field history can be controlled separately. This makes it useful for clarifying which aspects of artificial spin-ice dynamics arise from dipolar coupling itself and which require material-specific switching mechanisms.\par
\vspace{6pt}
[1] S.~H.~Skj{\ae}rv{\o} \textit{et al.}, Nat. Rev. Phys. \textbf{2}, 13 (2020).\par
[2] J.~H.~Jensen \textit{et al.}, Nat. Commun. \textbf{15}, 964 (2024).\par
[3] R.~L.~Novak and J.~P.~Sinnecker, J. Magn. Magn. Mater. \textbf{272--276}, 1557 (2004).
\end{document}
